% Sequence diagram 2: failed quiz -> suggestion -> the lesson in another form (R2, R5, R9)
\begin{sequencediagram}
    \scriptsize
    \def\unitfactor{0.72}
    \newthread{k}{Uczeń (klient)}
    \newinst[2.6]{as}{\texttt{assessment}}
    \newinst[1.9]{ls}{\texttt{learning\_style}}
    \newinst[3.1]{db}{Baza danych}

    \begin{call}{k}{\texttt{GET /lessons/\{id\}/quiz}}{as}{5 z 8 pytań}
    \end{call}

    \begin{call}{k}{\texttt{POST .../quiz/attempts}}{as}{40\%, wyjaśnienia}
        \begin{callself}{as}{ocena: 2 z 5}{}
        \end{callself}
        \begin{call}{as}{\texttt{record\_quiz\_result}}{ls}{\texttt{QuizFeedback}}
            \begin{callself}{ls}{R2: V 45 $\to$ 35, A 47, K 18}{}
            \end{callself}
            \begin{callself}{ls}{R5: sugestia A (ankieta 40\%)}{}
            \end{callself}
            \begin{call}{ls}{wynik, profil, zdarzenia}{db}{commit}
            \end{call}
        \end{call}
    \end{call}

    \begin{call}{k}{\texttt{PUT .../study/form} (A)}{ls}{wersja A lekcji}
        \begin{call}{ls}{wybór formy, zdarzenie}{db}{commit}
        \end{call}
    \end{call}

    \begin{call}{k}{\texttt{GET /lessons/\{id\}/quiz}}{as}{nowy zestaw (R9)}
    \end{call}
\end{sequencediagram}
